\chapter{Il Corto Circuito: Analisi delle Criticità del Sistema Vigente}

\section*{2.0 Dal modello teorico alla dinamica delle riforme}

Le criticità emerse nel capitolo precedente non rimangono confinate alla dimensione astratta della teoria previdenziale, ma trovano una traduzione immediata nella storia istituzionale italiana. La logica del sistema a ripartizione chiarisce infatti che, quando il tasso di crescita della popolazione attiva rallenta e la produttività non è sufficiente a compensarne gli effetti, la sostenibilità del sistema dipende in misura crescente da due variabili decisive: l'età effettiva di uscita dal lavoro e l'adeguatezza del montante contributivo accumulato~\cite{Bosi2018, RosenGayerRapallini2023}. In altri termini, la tenuta finanziaria non è un fatto puramente contabile, ma il risultato di un equilibrio tra demografia, mercato del lavoro e regole di accesso alle prestazioni.

È precisamente in questo punto di intersezione che si collocano le riforme pensionistiche italiane degli ultimi decenni. Esse non vanno lette come una sequenza discontinua di correzioni tecniche, bensì come il tentativo di ricomporre una frattura aperta tra il principio attuariale introdotto con il sistema contributivo nozionale e la pressione politico-sociale verso canali di uscita anticipata dal mercato del lavoro~\cite{LavoceRiformaDini, MEFRGS2025}. La questione delle deroghe, delle quote e delle finestre di flessibilità non costituisce dunque un capitolo accessorio, ma il luogo in cui si manifesta con maggiore evidenza la tensione tra sostenibilità finanziaria, adeguatezza delle prestazioni e consenso distributivo.

Per questo motivo, l'analisi che segue prende avvio dalle condizioni macroeconomiche e demografiche che hanno progressivamente eroso la base contributiva del sistema, per poi esaminare come le riforme Dini e Fornero abbiano ristabilito una coerenza di fondo, pur lasciando aperta la questione dell'equità intergenerazionale. Solo a partire da questa ricostruzione è possibile comprendere perché le successive misure di anticipo pensionistico abbiano inciso non soltanto sui conti dell'INPS, ma anche sulla credibilità stessa del patto previdenziale tra generazioni~\cite{MEFRGS2025, Bosi2018}.
\section*{2.1 La pressione macroeconomica: demografia e sostenibilità}

Il sistema pensionistico a ripartizione costruito dall'Italia tra gli anni Cinquanta e Settanta reggeva sotto precise condizioni demografiche ed economiche: molti lavoratori giovani, poche pensioni da pagare, crescita sostenuta del PIL. Quando quelle condizioni hanno cominciato a venire meno, il meccanismo ha rivelato una fragilità strutturale che nessuna manutenzione ordinaria poteva correggere~\cite{Bosi2018}.

\subsection*{2.1.1 La transizione demografica e l'inverno delle nascite}

Il tasso di fecondità italiano è passato dai 2,7 figli per donna del dopoguerra all'1,25 del 2023, mentre la speranza di vita ha continuato a crescere: secondo le proiezioni MEF 2025, è previsto un incremento ulteriore di 4,5 anni per i maschi e 3,9 per le femmine entro il 2070~\cite{MEFRGS2025}. L'effetto combinato è un aumento radicale dell'indice di dipendenza degli anziani, salito dal 31,1\% nel 2010 al 37,8\% nel 2023, con proiezioni che lo portano oltre il 57\% nel 2040~\cite{RosenGayerRapallini2023}.

In un sistema a ripartizione, questa dinamica si traduce direttamente nella formula sviluppata nel capitolo precedente: una contrazione di \(n\) comprime l'importo degli assegni oppure costringe ad alzare l'aliquota contributiva \(c\), a parità di salario medio \(S\) e produttività \(m\). L'aritmetica non offre scappatoie~\cite{Bosi2018}.

\subsection*{2.1.2 L'equazione di equilibrio della spesa pensionistica}

La sostenibilità macroeconomica si legge decomponendo il rapporto tra spesa pensionistica e PIL~\cite{RosenGayerRapallini2023}:

\[
\frac{S_p}{\text{PIL}} = \frac{N_p}{N_l} \times \frac{S_p / N_p}{\text{PIL} / N_l}
\]

Il primo termine è l'indice di dipendenza previdenziale; il secondo misura la generosità relativa del sistema. Una scomposizione più articolata include esplicitamente il tasso di occupazione~\cite{Bosi2018}:

\[
\frac{S_p}{\text{PIL}} = \frac{N_p}{N_a} \times \frac{N_a}{\text{POP}} \times \frac{\text{POP}}{N_l} \times \frac{S_p / N_p}{\text{PIL} / N_l}
\]

Il terzo termine, \(\text{POP}/N_l\), è l'inverso del tasso di occupazione e rivela come la bassa partecipazione femminile italiana abbia storicamente aggravato il divario tra massa contributiva e prestazioni da erogare~\cite{RosenGayerRapallini2023}. Nell'Italia degli anni Novanta tutte e quattro le componenti spingevano nella direzione sbagliata: il rapporto \(S_p/\text{PIL}\) aveva già raggiunto il 13,3\% nel 1990 su una traiettoria non decrescente, rendendo inevitabile la stagione delle grandi riforme descritta nel Capitolo 1~\cite{Bosi2018}.

\section*{2.2 Gli effetti delle riforme: sostenibilità acquisita, adeguatezza sacrificata}

Le riforme Dini e Fornero hanno conseguito il loro obiettivo primario. Il complesso degli interventi dal 2004 in poi è atteso ridurre il rapporto spesa/PIL di circa 60 punti percentuali cumulati al 2060 rispetto allo scenario controfattuale senza riforme, con circa un terzo del risparmio attribuibile specificamente alla riforma Fornero~\cite{MEFRGS2025}. La spesa si attesta al 15,4\% del PIL nel 2024, raggiungerà un picco vicino al 17\% nel biennio 2042-2044 con l'esaurimento delle coorti del baby boom, per poi scendere strutturalmente al di sotto del 14\% nel 2070 grazie all'applicazione generalizzata del metodo NDC~\cite{MEFRGS2025}.

Questa traiettoria discendente è il prodotto diretto del meccanismo di aggiustamento automatico introdotto dalla riforma Dini e accelerato dalla Fornero: all'aumentare di \(e_L\), il coefficiente \(\alpha_L = 1/e_L\) si riduce, comprimendo la pensione annua a parità di montante. Per il periodo 2025-2026 il coefficiente di trasformazione a 67 anni vale 5,608\%; le proiezioni demografiche indicano che potrebbe scendere attorno al 5,2\% nel 2040, riflettendo esclusivamente l'allungamento atteso della speranza di vita~\cite{MEFRGS2025}. Il MEF ha stimato che la rimozione permanente del solo meccanismo di adeguamento automatico dei requisiti all'aspettativa di vita comporterebbe un incremento del rapporto debito/PIL di circa 30 punti percentuali al 2070~\cite{LavoceRiformaDini, MEFRGS2025}.

Il sistema ha quindi risolto il problema del \textit{come} finanziare le pensioni future. Il problema del \textit{se} quelle pensioni saranno sufficienti è rimasto aperto, e i dati del MEF ne quantificano già oggi la dimensione~\cite{MEFRGS2025}.

\section*{2.3 Le criticità residue}

\subsection*{2.3.1 Adeguatezza degli assegni e rischio di povertà strutturale}

Il tasso di sostituzione lordo, ossia il rapporto tra la prima rata di pensione e l'ultima retribuzione, è destinato a scendere significativamente per le generazioni che si pensionano interamente sotto il regime contributivo. Secondo il Rapporto MEF aggiornamento 2025, per un lavoratore dipendente privato il tasso di sostituzione lordo si ridurrà dal 73,6\% nel 2010 al 58,4\% nel 2070 nello scenario base. Per i lavoratori autonomi il calo è ben più severo: dal 72,1\% al 46,7\% nello stesso orizzonte temporale~\cite{MEFRGS2025}.

Queste proiezioni assumono carriere lavorative continue e complete. La realtà del mercato del lavoro italiano, caratterizzata da ingresso tardivo, contratti discontinui e lunghi periodi di bassa contribuzione, suggerisce che per molti lavoratori i tassi di sostituzione effettivi saranno inferiori a quelli teorici. L'analisi di sensibilità del MEF mostra che una carriera con anzianità contributiva inferiore di due anni rispetto allo scenario base comprime ulteriormente il tasso di sostituzione già al di sotto della soglia del 58,4\%~\cite{LavoceGiovani, MEFRGS2025}.

Chi entra tardi nel mercato del lavoro, chi alterna lavoro dipendente a collaborazioni autonome, chi subisce interruzioni per ragioni familiari accumula un montante contributivo strutturalmente inferiore rispetto a quanto servirebbe per mantenere un tenore di vita dignitoso in vecchiaia. Il sistema NDC trasferisce questo rischio interamente sull'individuo: promette proporzionalità, non adeguatezza. Senza correttivi redistributivi o senza un secondo pilastro sviluppato, questa logica produce pensionati poveri per design~\cite{Bosi2018, LavoceGiovani}.

L'antidoto previsto dal sistema è la previdenza complementare. Con la partecipazione ai fondi pensione, il tasso di sostituzione di un lavoratore dipendente nel 2070 sale dal 58,4\% al 66\%, mentre per gli autonomi l'incremento è dal 46,7\% al 55,2\%. Il secondo pilastro rimane tuttavia sottosviluppato in Italia rispetto ai principali partner europei: la partecipazione volontaria è inferiore alle aspettative, e le carriere discontinue rendono difficile cumulare una contribuzione complementare adeguata~\cite{MEFRGS2025}.

\subsection*{2.3.2 Il costo delle deroghe e la rottura della neutralità attuariale}

La logica del sistema NDC è coerente solo se i requisiti di accesso vengono rispettati. Ogni deroga alla riforma Fornero, che si tratti di Quota 100, Quota 102, Quota 103 o di qualsiasi altro meccanismo di uscita anticipata, rompe la neutralità attuariale del sistema: consente di percepire la pensione per più anni del previsto senza un corrispondente aumento del montante, scaricando il differenziale di costo sulla collettività. Ogni anno di anticipo pensionistico non coperto da contributi aggiuntivi equivale a un trasferimento implicito dai lavoratori attivi di domani ai pensionati di oggi~\cite{LavoceRiformaDini, Bosi2018}.

Come documentato su Lavoce.info, il costo delle deroghe non è marginale. Ogni finestra di uscita anticipata che si mantiene strutturalmente nel sistema aumenta il numero di anni di erogazione attesi per coorte e abbassa il valore medio del coefficiente \(\alpha_L\) applicato al montante, poiché i lavoratori escono a età più giovani quando il coefficiente è più basso. Il risultato è un doppio svantaggio: pensioni più basse per chi esce prima e un onere aggiuntivo sul bilancio INPS che si ripercuote sulle generazioni future~\cite{LavoceRiformaDini, MEFRGS2025}.

Il paradosso è che ogni misura presentata come tutela dei lavoratori maturi porta con sé un costo pagato, con gli interessi, dai lavoratori più giovani. Il patto intergenerazionale che fonda il sistema a ripartizione regge solo finché le regole che lo governano mantengono una coerenza attuariale di fondo. Le deroghe iterate nel tempo non producono solo un costo contabile: erodono la credibilità stessa del contratto tra generazioni, rendendo razionale per i lavoratori giovani dubitare che il sistema onori le promesse che sta accumulando~\cite{Bosi2018, MEFRGS2025}.